% problem-000377 5 晶体结构与晶胞
\section*{5. 晶体结构与晶胞}
\textit{下列为分问。}

\subsection*{5-1}
A是紫红⾊晶体，分⼦量925.23，抗磁性。它通过 $\mathrm { R h C l _ { 3 } } { \cdot } 3 \mathrm { H _ { 2 } O }$ 和过量三苯膦 $\left( \mathrm { P P h } _ { 3 } \right)$ 的⼄醇溶液回流制得。画出A的⽴体结构。

\subsection*{5-2}
A可能的催化机理如下图所⽰（图中16e表⽰中⼼原⼦周围总共有16个电⼦）：

（图：）

画出 D 的结构式。

\subsection*{5-3}
确定图中所有配合物的中⼼原⼦的氧化态。

\subsection*{5-4}
确定A、C、D和E的中⼼离⼦的杂化轨道类型。

\subsection*{5-5}
⽤配合物的价键理论推测C和E显顺磁性还是抗磁性，说明理由。
